% Appendix A (writer: setting). Proofs for sec:setting. Source: research/tracks/foundations/notes-final.md, §§2, 3, 5, 6, 7.
\section{Function semimeasures and SFI: proofs}\label{app:setting}

Notation is that of \cref{sec:setting}. Tapes are $\omega=\omega_1\omega_2\cdots$; a finite union of cylinders $[w]$ is identified with the corresponding finite union of half-open dyadic intervals of $[0,1)$, which changes nothing up to a $\Unif$-null set.

\subsection{Möbius characterisation}\label{app:setting:mobius}

\begin{proof}[Proof of \cref{prop:setting:mobius}]
\emph{Convergence.} For each $R\subseteq T$, $\sum_{E:\dom E=R}\nu(D\cup E)\le\nu(D)$, by the second consequence of \cref{def:setting:weak} applied to the weak semimeasure $E\mapsto\nu(D\cup E)$ on data sets with domain disjoint from $\dom D$. So $\sum_{E:\dom E\subseteq T}\nu(D\cup E)\le2^{|T|}\nu(D)$, and $\Mob_T\nu(D)$ converges absolutely.

\emph{Necessity.} Let $\nu=\nu_P$ and $A_t:=\{F(t)\ne\bot\}$ for $t\in T$. By inclusion--exclusion,
\[
P\Big(F\supseteq D,\ \textstyle\bigcap_{t\in T}A_t^c\Big)=\sum_{R\subseteq T}(-1)^{|R|}P\Big(F\supseteq D,\ \textstyle\bigcap_{t\in R}A_t\Big),
\]
and $P(F\supseteq D,\bigcap_{t\in R}A_t)=\sum_{E:\dom E=R}\nu(D\cup E)$, a countable disjoint union over the labellings of $R$. Collecting terms gives $\Mob_T\nu(D)$, which is therefore a probability.

\emph{Sufficiency.} For finite $S\subseteq X$ and $v\in\Ybot^S$ let $v^{\rm def}$ be the data set of the defined entries of $v$ and $S_\bot(v)$ the set where $v=\bot$, and put $q_S(v):=\Mob_{S_\bot(v)}\nu(v^{\rm def})\ge0$. For $x\notin S$, split the sum defining $q_{S\cup\{x\}}(v,\bot)$ by whether $x\in\dom E$; absolute convergence allows the rearrangement, and
\[
q_{S\cup\{x\}}(v,\bot)=q_S(v)-\textstyle\sum_{a\in Y}q_{S\cup\{x\}}(v,a).
\]
So $\sum_{a\in\Ybot}q_{S\cup\{x\}}(v,a)=q_S(v)$, and $q_\emptyset=\nu(\emptyset)=1$. By induction each $q_S$ is a probability distribution on the countable set $\Ybot^S$, and the family is consistent under marginalisation. Kolmogorov's extension theorem gives a law $P$ on $\Ybot^X$ with these marginals, and $P(F\supseteq D)=q_{\dom D}(D)=\Mob_\emptyset\nu(D)=\nu(D)$. The cylinder probabilities of any such $P$ are the Möbius sums, by necessity; so $P$ is unique.
\end{proof}

\subsection{Machine laws and filling-in processes}\label{app:setting:machines}

\begin{proof}[Proof of \cref{prop:setting:machinelaw}]
Each coordinate of $f^T_\omega$ is determined by countably many cylinders, so $\omega\mapsto f^T_\omega$ is measurable into $\Ybot^X$. The set $\{\omega:f^T_\omega(x)=y\}=\bigcup_{(x,w,y)\in T}[w]$ is effectively open, uniformly in $(x,y)$ and an index of $T$; a finite intersection of effectively open sets is effectively open uniformly; and the measure of an effectively open set is lsc uniformly in its description. $\{f_\omega\supseteq\emptyset\}$ is the whole space.

\emph{Enumeration.} Let $(W_i)$ enumerate all c.e.\ sets of triples. Let $T_i$ be obtained by enumerating $W_i$ and accepting each triple unless it conflicts (same input, comparable strings, different labels) with a triple already accepted. Each $T_i$ is c.e.\ and consistent, and $T_i=W_i$ whenever $W_i$ is consistent. So $(T_i)$ lists exactly the function machines.
\end{proof}

\begin{remark}[Computable filling-in processes]\status{proved}\src{foundations Rem 2.7}\label{rem:setting:fillin}
A \emph{computable filling-in process} is a Markov chain $F_0=\emptyset\subseteq F_1\subseteq F_2\subseteq\cdots$ of finite partial functions whose transition kernel $\Prob[F_{t+1}=g\mid F_t=f]$ is supported on finitely many extensions $g\supseteq f$, given as an explicit list with rational probabilities computable from $(t,f)$. The law of its limit $F_\infty:=\bigcup_tF_t$ is a machine law, with an index computable from the kernel. Conversely, every machine law is the law of the limit of the nested process $F_t:=$ ``the part of $f_\omega$ certified by triples enumerated within $t$ steps with $|w|\le t$'', which depends on $\omega_{1:t}$; this process is Markov in $\omega_{1:t}$, not in $F_t$. So machine laws are exactly the limits of computable nested processes driven by a hidden finite state.
\end{remark}

\begin{proof}
The machine reads the tape bit by bit and samples $F_1,F_2,\dots$ in turn: to sample from a finitely supported law with rational probabilities it reads bits until the dyadic interval named by the bits read so far lies inside one of the consecutive half-open intervals whose lengths are those probabilities, and takes that outcome. This fails to terminate only when the tape's real is an endpoint, a null set. On query $x$ it runs the process until the first $t$ with $x\in\dom F_t$ and outputs $F_t(x)$; it diverges if there is none. So $T$ contains $(x,w,y)$ when the process, run on the bits of $w$, produces $F_1,\dots,F_t$ with $t$ the first index such that $x\in\dom F_t$, and $F_t(x)=y$. The process reads its bits in order, so on comparable strings it produces the same initial segment of the chain; with $F_t\subseteq F_{t+1}$ this makes $T$ consistent. Off the null set the sampled chain has the law of the process, and $f_\omega=F_\infty$. The converse is a restatement of \cref{def:setting:machine}.
\end{proof}

\subsection{Realising lower-semicomputable tree semimeasures}\label{app:setting:realise}

\begin{lemma}[Realisation]\status{proved (binary case known: Zvonkin--Levin 1970, unverified)}\src{foundations Lemma 3.1; check r2}\label{lem:setting:realise}
Let $A$ be a countable alphabet with a fixed enumeration and $\mu$ an lsc semimeasure on $A^*$: $\mu(\varepsilon)\le1$ and $\mu(s)\ge\sum_{a\in A}\mu(sa)$. There are finite unions of cylinders $I_t(s)$, increasing in $t$ and computable from $(t,s)$ and an index of $\mu$, such that $I(s):=\bigcup_tI_t(s)$ satisfies $I(sa)\subseteq I(s)$, $I(sa)\cap I(sb)=\emptyset$ for $a\ne b$, and $\Unif(I(s))=\mu(s)$ for every $s$; the cylinders enumerated for strings of one length are pairwise disjoint. Consequently $\omega\mapsto$ ``the longest $s$ with $\omega\in I(s)$'' is a monotone machine whose output extends $s$ with probability exactly $\mu(s)$, and the event ``the output extends $s$'' is certified by a finite prefix of $\omega$.
\end{lemma}

\begin{proof}
For binary sequences this is due to \citet{zvonkin1970complexity} (not checked against the source); the proof below covers countable alphabets. Let $\mu_t$ be computable dyadic lower approximations, nondecreasing in $t$, with $\mu_t\to\mu$ and $\mu_t(s)=0$ unless $|s|\le t$ and $s$ uses only the first $t$ letters (round a rational approximation down to a multiple of $2^{-t}$; this keeps monotonicity in $t$). On this finite tree define $\mu'_t$ from the leaves up by $\mu'_t(s):=\max\{\mu_t(s),\ \sum_a\mu'_t(sa)\}$. Then:
\begin{itemize}
\item $\mu'_t\le\mu$, by induction from the leaves, since $\sum_a\mu'_t(sa)\le\sum_a\mu(sa)\le\mu(s)$;
\item $\mu'_t$ is nondecreasing in $t$ (by induction from the leaves of the tree of stage $t$, using $\mu_{t+1}\ge\mu_t$), and $\mu'_t\ge\mu_t\to\mu$, so $\mu'_t\to\mu$;
\item $\mu'_t(s)\ge\sum_a\mu'_t(sa)$, and $\mu'_t(\varepsilon)\le\mu(\varepsilon)\le1$.
\end{itemize}
Build $I_t$ from the root down. Let $I_t(\varepsilon)\supseteq I_{t-1}(\varepsilon)$ be a finite union of cylinders of measure $\mu'_t(\varepsilon)$. Given $I_t(s)\supseteq I_{t-1}(s)$ of measure $\mu'_t(s)$, the room $I_t(s)\setminus\bigcup_aI_{t-1}(sa)$ has measure $\mu'_t(s)-\sum_a\mu'_{t-1}(sa)\ge\sum_a\big(\mu'_t(sa)-\mu'_{t-1}(sa)\big)$. All quantities are dyadic, so after refining to a common cylinder length we can add to each $I_{t-1}(sa)$ fresh cylinders inside the room, pairwise disjoint and of total measure $\mu'_t(sa)-\mu'_{t-1}(sa)$. The invariants are preserved, and $\Unif(I(s))=\lim_t\mu'_t(s)=\mu(s)$. Siblings stay disjoint because the $I_t$ increase in $t$. Each $I_t(s)$ is a finite union of cylinders, so $\omega\in I(s)$ is witnessed by a finite prefix.
\end{proof}

\subsection{What machine laws include}\label{app:setting:includes}

\begin{proof}[Proof of \cref{prop:setting:includes}]
(a) Apply \cref{lem:setting:realise} with alphabet $Y$, depth one and $\mu(\varepsilon):=1$, $\mu(y):=P(y\mid x)$, uniformly in $x$. This gives c.e.\ unions of cylinders $G_{x,y}$, disjoint in $y$, with $\Unif(G_{x,y})=P(y\mid x)$. The private stream of $x$ is $\omega^{(x)}_i:=\omega_{\langle j(x),i\rangle}$. Let $T$ contain $(x,w,y)$ when $w$ contains the first $L$ bits $u$ of $\omega^{(x)}$, for some $L$, and $[u]\subseteq G_{x,y}$ is certified. If $(x,w,y),(x,w',y')\in T$ with $w\sqsubseteq w'$, the two stream prefixes $u,u'$ are comparable, so $[u]\cap[u']\ne\emptyset$ and $G_{x,y}\cap G_{x,y'}\ne\emptyset$; hence $y=y'$. For distinct inputs the streams use disjoint tape coordinates, so they are independent and uniform, and $\nu_T(D)=\prod_{(x,y)\in D}\Unif(G_{x,y})=\nu_P(D)$. The index of $T$ is computable from one of $P$. For the claim on samplers: the output probabilities of a probabilistic machine are lsc (enumerate halting computations); the construction just given is the converse.

(b) $T:=\{(x,\varepsilon,\varphi(x)):x\in\dom\varphi\}$.

(c) Enumerate $X$ as $x_1,x_2,\dots$ and $\Ybot$ with $\bot$ last. Build nested partitions of $[0,1)$: the interval $J(v)$ of $v\in\Ybot^k$ is split, from left to right, into the intervals $J(va)$ of lengths $q_{k+1}(va)$, $a\in Y$ in order, and then $J(v\bot)$. Each left endpoint is $\mathrm{left}(v)+\sum_{a'<a}q_{k+1}(va')$, and $J(v\bot)$ starts at $\mathrm{left}(v)+q_k(v)-q_{k+1}(v\bot)$; all are computable reals. Let $T$ contain $(x_k,w,y)$ when the dyadic interval of $w$ is certified, by rational approximations of the endpoints, to lie in the interior of some $J(v)$ with $v\in\Ybot^k$ and $v_k=y\ne\bot$. Intervals of one level are disjoint and the levels are nested, so $T$ is consistent. Off the null set of endpoints a tape lies in the interior of exactly one $J(v)$ at each level, so $\Unif\{f_\omega(x_i)=v_i,\ i\le k\}=\Unif(J(v))=q_k(v)$ for every $v$, and the law of $f_\omega$ is $P$. If $Y$ is finite, each $q_S(v)=\Mob_{S_\bot(v)}\nu_P(v^{\rm def})$ is a finite alternating sum of computable reals. For a random total function $q_S(v)=\nu_P(v)$ if $v$ has no $\bot$, and $0$ otherwise.

(d) Realise $w$ by \cref{lem:setting:realise} (depth one, alphabet $\mathbb N$). This enumerates pairwise disjoint cylinders $[u]$, each labelled by some $i$, those labelled $i$ having total measure $w_i$; disjointness makes the set of strings $u$ prefix-free. Let $T:=\{(x,uw',y):[u]\text{ enumerated with label }i,\ (x,w',y)\in T_i\}$. As in \cref{thm:setting:universal}, $T$ is consistent, the tape after $u$ is uniform and independent of $u$, and $\nu_T(D)=\sum_iw_i\nu_{T_i}(D)$ for $D\ne\emptyset$; tapes outside all enumerated cylinders give the empty function.
\end{proof}

\subsection{Weak universality and its gap}\label{app:setting:weak}

\begin{proof}[Proof of \cref{thm:setting:weakuniversal}]
Let $(\phi_i)$ enumerate all lsc maps from data sets to $[0,1]$, with computable approximations $\phi_{i,t}$ nondecreasing in $t$. At stage $t$ let $U_t$ be the finite set of data sets whose inputs and labels are among the first $t$ inputs and labels. Define $\phi'_{i,t}$ on $U_t$ (and $0$ outside) by decreasing $|\dom D|$:
\[
\phi'_{i,t}(D):=\max\Big\{\phi_{i,t}(D),\ \max_{x\le t,\ x\notin\dom D}\ \textstyle\sum_{y\le t}\phi'_{i,t}(D+(x,y))\Big\}.
\]
If $\phi'_{i,t}(\emptyset)>1$, freeze: $\phi'_{i,s}:=\phi'_{i,t-1}$ for all $s\ge t$. Then $\phi'_{i,t}$ is nondecreasing in $t$ (by induction from the largest domains) and satisfies \eqref{eq:setting:weak} everywhere: inside $U_t$ by construction; outside $U_t$ both sides vanish, since a data set outside $U_t$ has no extension inside it; a frozen $\phi'_{i,t-1}$ is weak for the same reason. Let $\nu_i:=\lim_t\phi'_{i,t}$. For $x\notin\dom D$, monotone convergence gives $\sum_y\nu_i(D+(x,y))=\lim_t\sum_y\phi'_{i,t}(D+(x,y))\le\lim_t\phi'_{i,t}(D)=\nu_i(D)$, and $\nu_i(\emptyset)\le1$; so $\nu_i$ is an lsc weak semimeasure. If $\phi_i$ is already weak, induction from the largest domains gives $\phi'_{i,t}\le\phi_i$; so there is no freezing, and $\phi_{i,t}\le\phi'_{i,t}\le\phi_i$ gives $\nu_i=\phi_i$. Put $\Mw:=\sum_i\mm(i)\nu_i$. The weak inequalities are linear and $\sum_i\mm(i)\le1$, so $\Mw$ is an lsc weak semimeasure, and it dominates each $\nu_i$ with factor $\mm(i)\ge2^{-K(i)}$.
\end{proof}

\begin{proof}[Proof of \cref{thm:setting:weakvsmachine}]
(i) $\MF$ is an lsc weak semimeasure (\cref{prop:setting:machinelaw}); apply \cref{thm:setting:weakuniversal}.

(ii) Fix $S=\{s_1<\dots<s_k\}$ and put $\mu_S(y_1\cdots y_j):=\nu(\{(s_1,y_1),\dots,(s_j,y_j)\})$ for $j\le k$ and $\mu_S:=0$ on longer strings. The weak inequality at $s_{j+1}$ is exactly the semimeasure inequality for $\mu_S$, and $\mu_S$ is lsc uniformly in $(\nu,S)$. Let $I(\cdot)$ realise $\mu_S$ (\cref{lem:setting:realise}), and let $T_{\nu,S}$ contain $(s_j,w,y)$ when $[w]$ lies in an enumerated cylinder of some $I(u)$ with $|u|=j$ and $u_j=y$. The sets $I(u)$ with $|u|=j$ are pairwise disjoint, so $T_{\nu,S}$ is consistent; and a tape lies in $I(u^{(j)})$ for each $j\le k$ exactly when it lies in $I(u^{(k)})$ and $u^{(j)}$ is the length-$j$ prefix of $u^{(k)}$. So for $\dom D=S$, $\nu_{T_{\nu,S}}(D)=\mu_S(y_{s_1}\cdots y_{s_k})=\nu(D)$. (For proper subsets of $S$ only ``$\le$'' holds; check c6 confirms both facts on examples.) The index of $T_{\nu,S}$ is computable from $(\nu,S)$, so \cref{thm:setting:universal} gives $\MF(D)\ge2^{-K(\nu,S)-O(1)}\nu(D)$, and $K(\nu,S)\le K(\nu)+K(S)+O(1)$. Take $\nu=\Mw$, of complexity $O(1)$, for the second claim.

(iii) Use $\nutri$ of \cref{ex:setting:weaknotrpf} on the blocks $B_j:=\{3j,3j+1,3j+2\}$, $j<k$ (inputs as natural numbers). Put $\nutri^{\otimes k}(D):=\prod_{j<k}\nutri(D|_{B_j})$ if $\dom D\subseteq\bigcup_{j<k}B_j$, and $0$ otherwise. A one-point inequality at $x\in B_j$ involves only the $j$-th factor, and at $x$ outside the blocks both sides vanish; so $\nutri^{\otimes k}$ is weak (c1 checks $k=2$ exhaustively, 1458 inequalities; r1 also $k=3$). $\rho:=\sum_k\mm(k)\nutri^{\otimes k}$ is an lsc weak semimeasure, so $\Mw\ge2^{-O(1)}\rho$. For $\vec\imath\in\{1,2,3\}^k$ let $E_{\vec\imath}$ be the union over $j$ of the copy of $E_{i_j}$ in block $j$. These $3^k$ data sets are pairwise inconsistent, each has $|\dom E_{\vec\imath}|=2k$ and $\nutri^{\otimes k}(E_{\vec\imath})=2^{-k}$. $\MF$ is the containment functional of a random partial function, so $\sum_{\vec\imath}\MF(E_{\vec\imath})\le1$, and some $D_k:=E_{\vec\imath}$ has $\MF(D_k)\le3^{-k}$. But $\Mw(D_k)\ge2^{-O(1)}\mm(k)2^{-k}\ge2^{-k-K(k)-O(1)}$.

(iv) Each $-\log\Mw(E_{\vec\imath})\le k+K(k)+O(1)$ as in (iii). By concavity of $\log$, $3^{-k}\sum_{\vec\imath}-\log\MF(E_{\vec\imath})\ge-\log\big(3^{-k}\sum_{\vec\imath}\MF(E_{\vec\imath})\big)\ge k\log3$. Subtract.

(v) The argument of (iii) shows $\min_{\vec\imath}P(F\supseteq E_{\vec\imath})\le3^{-k}$ for every law $P$ of a random partial function, so the ratio to $\nutri^{\otimes k}(E_{\vec\imath})=2^{-k}$ is at most $(2/3)^k$. Let $P_1$ be the uniform law on the three total functions of \cref{ex:setting:weaknotrpf}(b); it satisfies $P_1\ge\frac23\nutri$ on all ten data sets of the support of $\nutri$ (checked by hand and by c9(b)). Its $k$-fold product over independent blocks has $P_1^{\otimes k}(F\supseteq D)=\prod_jP_1(F|_{B_j}\supseteq D|_{B_j})\ge(2/3)^k\nutri^{\otimes k}(D)$, with equality on the $E_{\vec\imath}$ (c9(b) for $k\le4$; r1).
\end{proof}

On the family of (iii), $\dom E_{\vec\imath}$ is determined by $(\vec\imath,k)$, so $K(\dom D_k)\le k\log3+K(k)+O(1)$. Larger label sets might give gaps closer to $K(\dom D)$; the foundations track did not work this out.

\subsection{The open question on random-function laws}\label{app:setting:rpf}

\paragraph{Reformulation (proof sketch; foundations Problem 3.6).} Suppose an lsc $\nu_P$ has approximations $\nu_{P_t}\uparrow\nu_P$, where each $P_t$ is a law of a random partial function on a finite domain with finitely many labels and rational probabilities, computable from $t$, and consecutive laws admit a coupling with $F_t\subseteq F_{t+1}$. The couplings form a polytope given by rational linear constraints, so a rational coupling can be found by search, and \cref{rem:setting:fillin} then realises $\nu_P$ exactly as a machine law. By Strassen's theorem for finite partially ordered sets \citep{strassen1965existence} (quoted, not checked against the source), such a coupling exists if and only if $P_t(U)\le P_{t+1}(U)$ for every up-set $U$ of the extension order. The containment events $\{F\supseteq D\}$ are only the principal up-sets, so $\nu_{P_t}\le\nu_{P_{t+1}}$ does not give this. The question becomes whether every lsc $\nu_P$ admits, at a constant-factor loss, approximations that increase on all up-sets.

Two steps are not written out: the Strassen step is quoted, and the reduction from general lower approximations of $\nu_P$ to finite rational laws $P_t$ is not carried out. An adversary argument in the spirit of \cref{ex:setting:weaknotrpf}, forcing repeated couplings on fresh blocks, was tried by the foundations track and settled neither direction.

\subsection{The normalised convention}\label{app:setting:normalised}

\begin{proposition}[Two-point criterion]\status{proved; computed}\src{foundations Prop 5.4 (= its referee's Prop B); checks c10(a), r13}\label{prop:setting:twopoint}
Let $P$ be the law of a random partial function on inputs $a,b$, with $A:=\{F(a)\ne\bot\}$, $B:=\{F(b)\ne\bot\}$ and $P(A),P(B)>0$. Consider
\[
(\star)\qquad P(B\mid F(a)=\alpha)=P(B\mid A)\ \text{ and }\ P(A\mid F(b)=\beta)=P(A\mid B)\quad\text{for all $\alpha,\beta$ of positive probability.}
\]
(i) If $(\star)$ holds, the normalised probability of $\{(a,\alpha),(b,\beta)\}$ is the same in both reveal orders, for every $(\alpha,\beta)$. (ii) Conversely, if every two-point probability $P(F(a)=\alpha,F(b)=\beta)$, $\alpha,\beta\in Y$, is positive (as for $\MF$) and the two orders agree for every $(\alpha,\beta)$, then $(\star)$ holds. (iii) Without positivity (ii) can fail.
\end{proposition}

\begin{proof}
Put $p(\alpha,\beta):=P(F(a)=\alpha,F(b)=\beta)$, $u(\alpha):=P(B\mid F(a)=\alpha)/P(B)$ and $v(\beta):=P(A\mid F(b)=\beta)/P(A)$. Revealing $a$ first gives the normalised product
\[
\frac{P(F(a)=\alpha)}{P(A)}\cdot\frac{p(\alpha,\beta)}{P(F(a)=\alpha,\,B)}=\frac{p(\alpha,\beta)}{P(A)P(B)\,u(\alpha)},
\]
and revealing $b$ first gives $p(\alpha,\beta)/(P(A)P(B)\,v(\beta))$. When $p(\alpha,\beta)>0$ the two agree exactly when $u(\alpha)=v(\beta)$; when $p(\alpha,\beta)=0$ both are $0$. (i) $(\star)$ gives $u\equiv P(A\cap B)/(P(A)P(B))\equiv v$. (ii) Agreement on all pairs forces $u(\alpha)=v(\beta)$ for all $\alpha,\beta$, so $u\equiv v\equiv\gamma$ for a constant $\gamma$. Averaging $u$ over $\alpha$ with weights $P(F(a)=\alpha)/P(A)$ gives $\gamma=P(A\cap B)/(P(A)P(B))$, which is $(\star)$. (iii) The law with $(F(a),F(b))$ equal to $(0,0)$ with probability $\frac14$, $(1,1)$, $(1,\bot)$, $(\bot,1)$ with probability $\frac18$ each, and $(\bot,\bot)$ with probability $\frac38$ has $P(B\mid F(a)=0)=1$, $P(B\mid F(a)=1)=\frac12$, $P(B\mid A)=\frac34$; it puts positive mass only on the diagonal defined pairs and is symmetric in $a,b$, so its normalised probabilities are order-independent (c10(a4)).
\end{proof}

The law of \cref{prop:setting:normalise}(b) violates $(\star)$: $P(B\mid F(a)=0)=\frac12\ne\frac23=P(B\mid A)$. Total laws always satisfy it. Check c10 found agreement between $(\star)$ and order independence on 400 of 400 random laws with positive two-point cells (r13: 300 of 300).

\begin{proof}[Proof of \cref{prop:setting:transfer}]
(a) Each normalised factor is the raw factor divided by $Z_t$. At the first step $Z_1=\MF\{f_\omega(x_1)\ne\bot\}\ge2^{-c}$, because a constant function is a machine law of complexity $O(1)$; combine with \cref{thm:setting:coding}(c).

(b) Each normalised conditional is a ratio of two positive lsc reals, hence computable relative to $\emptyset'$, uniformly. For a fixed input list, $y_{1:n}\mapsto\prod_t\MFbar(y_t\mid D_{t-1},x_t)$ is a probability distribution on $Y^n$: sum out $y_n$, then $y_{n-1}$, and so on; each factor sums to $1$, and a repeated input contributes a point mass on its stored label (c10(b) checks this exactly on 16,200 cases, repeated inputs included). Apply the conditional coding theorem relative to $\emptyset'$.

(c) The normalised predictor is a history-aware predictor issuing probability vectors. At each step, conditionally on $(h_{t-1},x_t)$, Gibbs' inequality gives $\Ex[-\log\MFbar(y_t\mid D_{t-1},x_t)\mid h_{t-1},x_t]\ge H(y_t\mid h_{t-1},x_t)$; take expectations and sum.
\end{proof}

\subsection{Invariance, coding bounds and computability}\label{app:setting:basic}

\begin{proof}[Proof of \cref{thm:setting:invariance}]
(a) Each of $\MF,M'_F$ is a machine law in the other's enumeration, so each dominates the other up to a constant factor by \cref{thm:setting:universal}; the same argument applies to any two function machines whose laws dominate all machine laws up to constants.

(b) $\alpha^{-1}$ and $\beta^{-1}$ are partial computable (by search). The set of triples $(\alpha(x),w,\beta(y))$ with $(x,w,y)\in T_U$ is a function machine over $X',Y'$ (consistent because $\alpha$ is injective), whose random function is $\beta\circ f_\omega\circ\alpha^{-1}$, undefined off the range of $\alpha$. Its law $\nu'$ satisfies $\nu'(D^{\alpha\beta})=\Unif\{\beta(f_\omega(x))=\beta(y)\ \forall(x,y)\in D\}=\MF(D)$, since $\beta$ is injective. Its index is computable from indices of $\alpha,\beta$, so $M'_F(D^{\alpha\beta})\ge2^{-K(\alpha,\beta)-O(1)}\MF(D)$. Conversely, the machine with random function $\beta^{-1}\circ f'_\omega\circ\alpha$ on $X$ has law $\nu$ with $\nu(D)=M'_F(D^{\alpha\beta})$, so $\MF(D)\ge2^{-K(\alpha,\beta)-O(1)}M'_F(D^{\alpha\beta})$.

(c) For $W$: a program for $\varphi$ becomes a program for $\beta\circ\varphi\circ\alpha^{-1}$ by a prefix of length $K(\alpha,\beta)+O(1)$, and conversely. For $\Mprod$: transform each conditional semimeasure in the same way. For $\Mw$: given a weak $\nu$ on $X\times Y$, put $\nu'(D'):=\nu(D)$ if $D'=D^{\alpha\beta}$ and $\nu'(D'):=0$ otherwise; $\nu'$ is lsc (search for preimages). If $D'=D^{\alpha\beta}$ and $x'\notin\dom D'$, the extensions of $D'$ at $x'$ of the form $E^{\alpha\beta}$ have $x'=\alpha(x)$ and labels $\beta(y)$, so the one-point inequality transfers; if $D'$ is not an image, neither is any extension of it. The converse transformation is symmetric.
\end{proof}

\begin{proof}[Proof of \cref{thm:setting:coding}]
(a) $P(E\mid S):=\MF(E)$ for $\dom E=S$ is lsc uniformly in $(E,S)$ and sums to at most $1$ over $E$ (second consequence of \cref{def:setting:weak}). The conditional coding theorem gives $K(E\mid S)\le-\log\MF(E)+O(1)$. Take $E=D|_T$, $S=T$, and use $\MF(D)\le\MF(D|_T)$.

(b) First term: the shortest consistent program $p$ gives the machine law $\ind[\varphi_{V(p)}\supseteq D]$, whose complexity is at most $|p|+O(1)$, so $\MF(D)\ge2^{-\KF(D)-O(1)}$ (equivalently, $\MF\ge2^{-O(1)}W$ by \cref{thm:versions:chain}). Second term: $\mm(y\mid x):=\sum_{q:V(q,x)=y}2^{-|q|}$, with $V(\cdot,x)$ the conditional universal prefix machine, is an lsc conditional semimeasure with $\mm(y\mid x)\ge2^{-K(y\mid x)}$; its product law is a machine law of complexity $O(1)$ (\cref{prop:setting:includes}(a)). Last inequality: $\KF(D)\le K(D)+O(1)$ (a program computes the lookup table of $D$ from a shortest description of $D$) and $K(D)\le K(\dom D)+K(D\mid\dom D)+O(1)$.

(c) Combine (a) with $T=\{x\}$ and (b).
\end{proof}

\begin{proof}[Proof of \cref{prop:setting:computability}]
(a) \Cref{prop:setting:machinelaw} and the remark after \cref{def:setting:sfi}.

(c) Take $X=\mathbb N$ and suppose the raw conditional is computable. Define $f$ by recursion: with $D_t:=\{(s,f(s)):s<t\}$, let $f(t)$ be the first $y\in\{0,1\}$ for which approximation certifies $\MF(y\mid D_t,t)<0.6$. The two conditionals sum to at most $1$, so one of them is at most $\frac12$, and the search halts. $f$ is total computable, and $\MF(D_n)=\prod_{t<n}\MF(f(t)\mid D_t,t)<0.6^n$, which contradicts $\MF(D_n)\ge2^{-K(f)-O(1)}$ for large $n$. For the normalised conditional the same construction gives normalised loss at least $n\log(1/0.6)$, while it is at most $\LSFI(D_n)\le K(f)+O(1)$ (\cref{prop:setting:normalise}(a)).

(b) If $\MF$ were computable, the raw conditional would be a ratio of positive computable reals, hence computable.

(d) Lsc reals are computable relative to $\emptyset'$; the conditionals are ratios of positive lsc reals ($\MF(D)$ and, for the normalised one, $\Znorm(D,x)$).
\end{proof}
